\documentclass[10pt,journal,final,a4paper]{IEEEtran}

% --- 1. Mathematik & Symbole ---
\usepackage{amsmath} 
\usepackage{amssymb}
\usepackage{amsfonts}

% --- 2. Code-Blöcke (Zuerst laden wegen fvextra/lineno-Interaktionen) ---
% \usepackage{minted} 
\usepackage{listings}

% --- 3. Text & Sprache (csquotes MUSS nach minted/fvextra geladen werden) ---
\usepackage{csquotes} 

% --- 4. Tabellen & Grafiken ---
\usepackage{color,array}
\usepackage{graphicx}
\usepackage{multirow} 
\usepackage{booktabs} 
\usepackage{float}
\usepackage{siunitx}
\usepackage{tabularx}

% --- 5. IEEEtran Layout-Helfer (algorithmic gelöscht, da nicht vorhanden/benötigt) ---
\usepackage[caption=false,font=normalsize,labelfont=sf,textfont=sf]{subfig}
\usepackage{textcomp}
\usepackage{stfloats}
\usepackage{verbatim}
\usepackage{balance}

% --- 6. Links & Referenzen (Immer ganz ans Ende der Präambel) ---
\usepackage[colorlinks,urlcolor=blue,linkcolor=blue,citecolor=blue]{hyperref}
\usepackage{orcidlink} % benötigt hyperref, daher danach laden
\usepackage[capitalise,nameinlink,noabbrev]{cleveref} % MUSS als letztes geladen werden

% --- IEEEtran spezifische Definitionen ---
\hyphenation{op-tical net-works semi-conduc-tor IEEE-Xplore}
\def\BibTeX{{\rm B\kern-.05em{\sc i\kern-.025em b}\kern-.08em
    T\kern-.1667em\lower.7ex\hbox{E}\kern-.125emX}}

\newtheorem{theorem}{Theorem}
\newtheorem{lemma}{Lemma}
\setcounter{page}{1}

\begin{document}

\title{GoTogether} 

% TODO FOR PUBLICATION: edit author list and affiliations
\author{Simon Völkl\,\orcidlink{0009-0005-5334-9673}, 
        Johannes Schieder\,\orcidlink{0009-0004-0096-3808}%
\thanks{The authors are with the Department of Electrical Engineering, Media, and Computer Science, Ostbayerische Technische Hochschule Amberg-Weiden, 92224 Amberg, Germany (e-mail: \{s.voelkl2, j.schieder\}@oth-aw.de).}%
}

% TODO FOR PUBLICATION: Update date, title
\markboth{TECHNICAL REPORT, JULY 2026}
{Völkl \MakeLowercase{\textit{et al.}}: GoTogether}

\maketitle

\begin{abstract}
    % \textbf{Background:} 
    % \textbf{Objective:} 
    % \textbf{Methods:} 
    % \textbf{Results:} 
    % \textbf{Conclusion:} 
\end{abstract}

\begin{IEEEkeywords}
    % A, B, C, D
\end{IEEEkeywords}


% \tableofcontents

\section{Introduction}
\label{sec:introduction}

% TODO FOR PUBLICATION: update introduction
\IEEEPARstart{E}{mbedded}...
The source code and documentation for the GoTogether project are publicly available on GitHub~\cite{gotogether2026github}.

% \begin{figure}[htb]
%     \centering
%     \includegraphics[width=0.6\linewidth]{../assets/logo/logo_cut.png}
%     \caption{GoTogether Project Logo. Designed by Simon Völkl.}
%     \label{fig:logo}
% \end{figure}

\subsection{Mission Statement}
\label{subsec:mission_statement}


\subsection{Document Structure}
\label{subsec:document_structure}


\section{Technical Concept}
\label{sec:technical_concept}

\subsection{High-Level Architecture}
\label{subsec:high_level_architecture}

\subsection{Components}
\label{subsec:components_concept}

\subsection{Project Management}
\label{subsec:project_management}


\subsection{Task Dependencies}
\label{subsec:task_dependencies}

\subsection{Task Distribution}
\label{subsec:task_distribution}

% To coordinate development and ensure progress, tasks were tracked and managed using a GitHub Project board integrated directly with the code repository. This allowed issues to be linked to commits and Pull Requests (PRs).

% To maintain code quality and structural integrity across the shared codebase, the development followed several guidelines:

% \begin{itemize}
%     \item \textbf{Branch Protection:} Direct commits to the \texttt{main} branch were restricted. All feature updates and bug fixes were developed on separate branches.
%     \item \textbf{Code Review:} Merging any branch into \texttt{main} required a Pull Request accompanied by the other team member's review and approval.
%     \item \textbf{CI Pipeline}: A GitHub Actions pipeline was configured to automatically verify the compilation of the ESP32 firmware for both test and production builds.
%     \item \textbf{Code Coverage Targets:} A unit test coverage of about 60\% was targeted for the core ESP32 modules, while the Android application aimed for basic unit tests. Node-RED flows were excluded from automated testing due to their visual nature.
%     \item \textbf{Docstrings:} All public methods and classes were documented with docstrings to enhance code readability and accessibility.
% \end{itemize}

% Transparency note: The unit test writing and docstring generation were partly assisted by AI.

\section{Discussion}
\label{sec:discussion}

\section{Summary and Future Work}
\label{sec:summary-future-work}

% --- References ---
% The \balance command ensures that the two columns on the last
% page end cleanly at the same height. Uncomment as soon as the text is finished!
\balance 
\bibliographystyle{IEEEtran}
\bibliography{references}

\end{document}